\chapter{Formy dwuliniowe i kwadratowe}
\label{ch:lin9-formy-kwadratowe}
\index{forma kwadratowa@Forma kwadratowa}
\index{forma dwuliniowa symetryczna@Forma dwuliniowa symetryczna}
\index{ortogonalność@Ortogonalność}
\index{radykał formy@Radykał formy}
\index{kongruencja macierzy@Kongruencja macierzy}

\noindent\textbf{Cel rozdziału.}
W Rozdziale~\ref{ch:lin8-wieloliniowe}
poznaliśmy formy dwuliniowe, ich macierze
i zasady zmiany bazy. Zbadamy teraz strukturę
form symetrycznych i odpowiadających im
form kwadratowych. Zdefiniujemy ortogonalność,
wyprowadzimy rozkłady ortogonalne,
udowodnimy twierdzenie o diagonalizacji
i prawo bezwładności Sylvestera.
Oddzielnie przeanalizujemy formy kwadratowe
w charakterystyce \(2\) i formy naprzemienne.
Rozdział nie korzysta z iloczynów tensorowych
ani z teorii wartości własnych.

Przez cały rozdział \(K\) oznacza ciało
(\(0_K\ne1_K\)); wszystkie przestrzenie
wektorowe są nad \(K\). W razie
potrzeby skończoność wymiaru,
charakterystykę i uporządkowanie
ciała wskazujemy jawnie.
Wymiar \(0\) jest dopuszczalny.

\section{Formy symetryczne i ortogonalność}
\label{sec:lin9-ortogonalnosc}

\begin{definicja}[Ortogonalność i radykał]
\label{def:lin9-ortogonalnosc}
Niech \(V\) będzie przestrzenią
wektorową, a
\(b\colon V\times V\to K\)
symetryczną formą dwuliniową.
Dla \(u,v\in V\) piszemy
\(u\perp_b v\), jeśli \(b(u,v)=0_K\).
Dla podprzestrzeni \(U\subseteq V\)
jej \emph{dopełnieniem ortogonalnym}
względem \(b\) nazywamy
\[
 U^{\perp_b}
 :=\{v\in V:\ b(v,u)=0_K
                   \text{ dla każdego }u\in U\}.
\]
\emph{Radykał} formy \(b\) to
\[
 \operatorname{Rad}(b):=V^{\perp_b}.
\]
Formę \(b\) nazywamy
\emph{niezdegenerowaną},
jeśli \(\operatorname{Rad}(b)=\{0_V\}\).
Ograniczenie formy do \(U\)
oznaczamy \(b|_{U\times U}\).
\end{definicja}

\begin{stwierdzenie}[Podstawowe własności ortogonalności]
\label{prop:lin9-perp-wlasnosci}
Niech \(V,b\) będą jak w
Definicji~\ref{def:lin9-ortogonalnosc}.
Dla każdej podprzestrzeni
\(U\subseteq V\) zbiór
\(U^{\perp_b}\) jest podprzestrzenią.
Ponadto
\[
 \operatorname{Rad}(b)\subseteq U^{\perp_b},
 \qquad
 U\subseteq (U^{\perp_b})^{\perp_b}.
\]
Dla \(U_1\subseteq U_2\)
zachodzi
\(U_2^{\perp_b}\subseteq U_1^{\perp_b}\).
Radykał ograniczenia \(b|_{U\times U}\)
jest równy \(U\cap U^{\perp_b}\).
\end{stwierdzenie}
\begin{proof}
Dla każdego \(u\in U\)
odwzorowanie \(v\mapsto b(v,u)\)
jest liniowe. Zatem
\(U^{\perp_b}\) jest częścią wspólną
jąder tych odwzorowań, a więc
podprzestrzenią.
Jeśli \(v\in\operatorname{Rad}(b)\),
jest ortogonalny do wszystkich
wektorów \(V\), w szczególności
do każdego \(u\in U\).
Daje to pierwsze zawieranie.
Jeśli \(u\in U\) oraz
\(v\in U^{\perp_b}\), to
\(b(v,u)=0_K\), a przez
symetrię również \(b(u,v)=0_K\).
Oznacza to \(u\in
(U^{\perp_b})^{\perp_b}\).
Zawieranie przy \(U_1\subseteq U_2\)
wynika wprost z liczby
warunków definiujących
dopełnienie ortogonalne.
Wektor \(u\in U\) leży
w radykale ograniczenia
dokładnie wtedy, gdy
\(b(u,u')=0_K\) dla każdego
\(u'\in U\), czyli gdy
\(u\in U^{\perp_b}\).
\end{proof}

\begin{twierdzenie}[Wymiar dopełnienia ortogonalnego]
\label{thm:lin9-perp-wymiar}
Niech \(\dim V=n<\infty\),
\(b\colon V\times V\to K\)
będzie symetryczną formą
dwuliniową, a \(U\subseteq V\)
podprzestrzenią.
Wtedy
\[
 \dim U^{\perp_b}
 =n-\dim U+\dim(U\cap\operatorname{Rad}(b)).
\]
W szczególności jeśli \(b\)
jest niezdegenerowana, to
\(\dim U^{\perp_b}=n-\dim U\)
oraz
\((U^{\perp_b})^{\perp_b}=U\).
\end{twierdzenie}
\begin{proof}
Niech \(R=\operatorname{Rad}(b)\).
Zdefiniujmy liniowe
\(L\colon V\to U^*\)
przez
\(L(v)(u)=b(v,u)\).
Jądro \(L\) jest dokładnie
\(U^{\perp_b}\).
Rozważmy też odwzorowanie
\(M\colon U\to V^*\)
dane przez \(M(u)(v)=b(u,v)\).
Jego jądro jest równe \(U\cap R\).
Ponieważ \(b\) jest symetryczna,
\(L\) jest odwzorowaniem dualnym
do \(M\) po utożsamieniu
\(V\cong V^{**}\) i \(U\cong U^{**}\).
Wszystkie przestrzenie są
skończenie wymiarowe, więc
z
Twierdzenia~\ref{thm:lin7-rzad-dualnego}
wynika
\(\operatorname{rk}L
=\operatorname{rk}M
=\dim U-\dim(U\cap R)\).
Twierdzenie o rzędzie i jądrze
daje
\[
 \dim U^{\perp_b}
 =\dim\ker L
 =n-\operatorname{rk}L
 =n-\dim U+\dim(U\cap R).
\]
Jeśli \(R=0\), otrzymujemy
\(\dim U^{\perp_b}=n-\dim U\).
Stosując ten sam wzór do
\(U^{\perp_b}\), dostajemy
\(\dim (U^{\perp_b})^{\perp_b}
=n-(n-\dim U)=\dim U\).
Z
Stwierdzenia~\ref{prop:lin9-perp-wlasnosci}
mamy zawieranie
\(U\subseteq (U^{\perp_b})^{\perp_b}\).
Równość wymiarów kończy dowód.
\end{proof}

\begin{twierdzenie}[Ortogonalna suma prosta]
\label{thm:lin9-suma-ortogonalna}
Niech \(V\) będzie skończenie
wymiarową przestrzenią,
\(b\colon V\times V\to K\)
symetryczną formą dwuliniową
i \(U\subseteq V\) podprzestrzenią.
Jeśli ograniczenie
\(b|_{U\times U}\)
jest niezdegenerowane, to
\[
 V=U\oplus U^{\perp_b}.
\]
Dla \(x=u+w\), gdzie
\(u\in U\), \(w\in U^{\perp_b}\),
zachodzi
\(b(x,x)=b(u,u)+b(w,w)\).
Nie wymagamy niezdegenerowania
formy na całej przestrzeni \(V\).
\end{twierdzenie}
\begin{proof}
Z
Stwierdzenia~\ref{prop:lin9-perp-wlasnosci}
wynika
\(U\cap U^{\perp_b}
=\operatorname{Rad}(b|_{U\times U})
=\{0\}\).
W szczególności
\(U\cap\operatorname{Rad}(b)=\{0\}\).
Zatem
Twierdzenie~\ref{thm:lin9-perp-wymiar}
daje
\(\dim U^{\perp_b}
=n-\dim U\).
W związku z tym
\(\dim(U+U^{\perp_b})=n\),
a suma jest bezpośrednia,
więc obejmuje całe \(V\).
Dla \(u\in U\),
\(w\in U^{\perp_b}\)
symetria daje
\(b(u,w)=b(w,u)=0_K\).
Rozwijając \(b(u+w,u+w)\),
otrzymujemy wskazany wzór.
\end{proof}

\begin{przyklad}[Dopełnienie ortogonalne w formie zdegenerowanej]
\label{prz:lin9-radykal-przyklad}
Niech \(V=\mathbb R^3\),
\(b(x,y)=x_1y_1+x_2y_2\).
Wtedy
\(\operatorname{Rad}(b)
=\operatorname{span}\{e_3\}\).
Dla \(U=\operatorname{span}\{e_1\}\)
mamy \(U^{\perp_b}
=\operatorname{span}\{e_2,e_3\}\).
Ograniczenie \(b\) do \(U\)
jest niezdegenerowane i rzeczywiście
\(V=U\oplus U^{\perp_b}\).

Natomiast dla
\(U_0=\operatorname{span}\{e_3\}\)
otrzymujemy
\(U_0^{\perp_b}=V\).
Stąd
\(U_0+U_0^{\perp_b}=V\),
ale suma nie jest bezpośrednia,
ponieważ \(e_3\) należy do
obu podprzestrzeni.
Wzór wymiarowy daje
\(\dim U_0^{\perp_b}
=3-1+1=3\).
\end{przyklad}

\section{Formy kwadratowe i polaryzacja}
\label{sec:lin9-kwadratowe}

\begin{definicja}[Forma kwadratowa]
\label{def:lin9-kwadratowa}
Niech \(V\) będzie przestrzenią
wektorową nad ciałem \(K\).
\emph{Formą kwadratową} na \(V\)
nazywamy funkcję \(q\colon V\to K\)
taką, że
\[
 q(av)=a^2q(v)
 \quad(a\in K,\ v\in V),
\]
a odwzorowanie
\[
 b_q(u,v):=q(u+v)-q(u)-q(v)
\]
jest dwuliniowe.
Formę \(b_q\) nazywamy
\emph{polaryzacją} \(q\).
Jest ona symetryczna,
ponieważ \(u+v=v+u\).
\end{definicja}

\begin{stwierdzenie}[Własności polaryzacji]
\label{prop:lin9-polaryzacja}
Dla każdej formy kwadratowej
\(q\colon V\to K\)
zachodzą
\[
 q(0_V)=0_K,\qquad
 q(-v)=q(v),\qquad
 b_q(v,v)=2_Kq(v).
\]
Ponadto
\[
 q(u+v)=q(u)+q(v)+b_q(u,v).
\]
Jeżeli \(2_K\ne0_K\), to
\[
 q(v)=\tfrac12 b_q(v,v),
 \qquad
 b_q(u,v)=\tfrac12
 \bigl(q(u+v)-q(u-v)\bigr).
\]
W tej charakterystyce
przyporządkowania
\(q\mapsto b_q\)
i \(b\mapsto[v\mapsto\tfrac12 b(v,v)]\)
są wzajemnie odwrotnymi
odwzorowaniami między
przestrzenią form kwadratowych
i przestrzenią symetrycznych
form dwuliniowych.
\end{stwierdzenie}
\begin{proof}
Z jednorodności, dla \(a=0_K\),
wynika \(q(0_V)=0_K\).
Dla \(a=-1_K\)
otrzymujemy \(q(-v)=q(v)\).
Wstawiając \(u=v\) do
definicji polaryzacji,
otrzymujemy
\[
 b_q(v,v)
 =q(2_Kv)-2_Kq(v)
 =(4_K-2_K)q(v)=2_Kq(v).
\]
Równość dla \(q(u+v)\)
jest przekształceniem
definicji polaryzacji.
Jeśli \(2_K\ne0_K\), możemy
podzielić przez \(2_K\).
Z dwuliniowości \(b_q\)
i symetrii wynika
\[
 q(u+v)-q(u-v)
 =b_q(u,v)-b_q(u,-v)
 =2_Kb_q(u,v).
\]
Jeśli odwrotnie
\(b\colon V\times V\to K\)
jest symetryczna, zadajmy
\(q_b(v)=\tfrac12 b(v,v)\).
Dwuliniowość daje
\(q_b(av)=a^2q_b(v)\)
oraz
\[
 b_{q_b}(u,v)
 =\tfrac12\bigl(
  b(u,v)+b(v,u)\bigr)=b(u,v).
\]
Dla istniejącego \(q\)
pierwszy wzór daje
\(q_{b_q}=q\).
Obie konstrukcje są więc
wzajemnie odwrotne.
\end{proof}

\begin{twierdzenie}[Współrzędne formy kwadratowej]
\label{thm:lin9-macierz-kwadratowej}
Niech \(\dim V=n<\infty\)
i \(B=(e_1,\ldots,e_n)\)
będzie uporządkowaną bazą \(V\).
Wtedy każda forma kwadratowa
\(q\colon V\to K\)
ma jednoznaczne współczynniki
\(a_i,c_{ij}\in K\),
\(1\leq i\leq n\),
\(1\leq i<j\leq n\),
takie, że dla
\(x=\sum_i x_ie_i\)
zachodzi
\[
 q(x)=\sum_{i=1}^{n}a_ix_i^2+
       \sum_{1\leq i<j\leq n}
                         c_{ij}x_ix_j.
\]
Są to
\[
 a_i=q(e_i),\qquad
 c_{ij}=b_q(e_i,e_j).
\]
Odwrotnie, każdy wybór
tych współczynników definiuje
formę kwadratową.
Jeżeli \(2_K\ne0_K\),
istnieje dokładnie jedna
macierz symetryczna \(A\)
spełniająca \(q(x)=X^{\mathsf T}AX\),
gdzie \(X=[x]_B\); jest nią
\[
 A_{ii}=a_i,\qquad
 A_{ij}=A_{ji}=\tfrac12c_{ij}
 \quad(i<j).
\]
Przy zmianie bazy z macierzą
przejścia \(P\) obowiązuje
\(A'=P^{\mathsf T}AP\).
\end{twierdzenie}
\begin{proof}
Ze wzoru
\(q(u+v)=q(u)+q(v)+b_q(u,v)\),
stosowanego indukcyjnie
do sumy \(\sum_i x_ie_i\),
otrzymujemy
\[
 q\!\left(\sum_i x_ie_i\right)
 =\sum_i q(x_ie_i)+
  \sum_{i<j}b_q(x_ie_i,x_je_j).
\]
Jednorodność i dwuliniowość
dają żądane rozwinięcie.
Po wstawieniu \(e_i\)
odczytujemy \(a_i=q(e_i)\),
a po wstawieniu \(e_i+e_j\)
wraz z wartościami na \(e_i,e_j\)
odczytujemy
\(c_{ij}=b_q(e_i,e_j)\).
To dowodzi jednoznaczności
nawet nad ciałem skończonym.

Dla dowolnych \(a_i,c_{ij}\)
prawa strona jest jednorodna
stopnia \(2\), a jej polaryzacja
jest równa
\[
 b_q(x,y)=\sum_i2_Ka_ix_iy_i+
 \sum_{i<j}c_{ij}(x_iy_j+x_jy_i),
\]
która jest dwuliniowa.
Zatem definiuje formę
kwadratową.
Jeśli \(2_K\ne0_K\), podana
macierz symetryczna daje
dokładnie to rozwinięcie
po wymnożeniu
\(X^{\mathsf T}AX\);
jej jednoznaczność wynika
ze wzorów na \(a_i,c_{ij}\).
Zmianę bazy otrzymujemy
z prawa zmiany bazy formy
dwuliniowej \(b_q\) i równości
\(q(x)=\tfrac12b_q(x,x)\),
por. Twierdzenie~\ref{thm:lin8-macierz-formy}.
Wymiar \(0\) odpowiada
pustym sumom i jedynej
zerowej formie kwadratowej.
\end{proof}

\begin{przyklad}[Polaryzacja w dodatniej charakterystyce]
\label{prz:lin9-char2-polar}
Niech \(K=\mathbb F_2\) i
\(V=K^2\). Określmy
\(q(x_1,x_2)=x_1x_2+x_1^2\).
Jest to forma kwadratowa,
a jej polaryzacja wynosi
\[
 b_q(x,y)=x_1y_2+x_2y_1.
\]
Nie zależy od składnika
\(x_1^2\), gdyż
\((x_1+y_1)^2-x_1^2-y_1^2=0\).
Co więcej, forma
\(q_0(x_1,x_2)=x_1x_2\)
ma dokładnie tę samą
polaryzację, lecz
\(q(e_1)=1\), \(q_0(e_1)=0\).
Dlatego w charakterystyce \(2\)
polaryzacja nie wyznacza
formy kwadratowej jednoznacznie.
\end{przyklad}
